% Chapter 1 -- the idea. Update when the overall flow (pipeline.py, companies.py, background/, suggest/,
% affinity/, registers/, applications/, mail/, export/, chat/, recruiter_sim/, jobsearch/) changes.
\chapter{What it does and the idea}

ReverseHire is a desktop application (Python, PySide6) with one goal: to go \emph{from
the CV to the companies that could want this profile}. It reads one CV, works out who you are,
finds the companies around you (or in a province, a region, the country), reads them, and says for
each how well the profile fits what it seems to need, with the reasons. It also simulates the
screening software an employer may use, searches postings and tracks the applications you send.

\section{``Reverse'' hiring}
The usual job search starts from the postings that happen to be online. ReverseHire first works
out \emph{who you are}, then looks for the \emph{companies that could employ such a profile}, whether
they publish postings or not, and reads their websites to see which are hiring and how well you fit. A company with no posting but a plausible need is still
shown, as a candidate for a \emph{spontaneous application}. A company you did not know can come
from the free regional register (chapter~\ref{ch:register}) or from your own assistant, whose
suggestions are checked before they are listed (chapter~\ref{ch:interest}).

\section{The flow in one table}
{\small
\begin{tabularx}{\textwidth}{@{}l l Y@{}}
\toprule
\hdr \thead{Step} & \thead{Code} & \thead{What happens} \\
\midrule
1. Read the CV     & \code{cv/}        & Text (local OCR for scans), sections, skills, languages; loaded when you choose or drop the file. \\
2. Build the profile & \code{cv/}, \code{geo/sectors.py} & Roles, skills, sectors, ATECO categories, shown in English; every edit saved at once. \\
3. Test the CV     & \code{recruiter\_sim/} & Simulated screening systems: parse, knock-outs, score in per cent, fixes. \\
4. Find companies  & \code{discovery/}, \code{registers/} & Local map data, regional register, Wikidata, favourites; ranked. \\
5. Read the sites  & \code{web/}       & Careers page, ATS boards, postings, legal identifiers, links. \\
6. Score           & \code{scoring/}, \code{affinity/} & The \emph{match}; then the \emph{affinity} and its verdict. \\
7. Who could be interested & \code{background/}, \code{suggest/} & Pre-screen of every known company, verified assistant suggestions, deep analysis when the PC is idle. \\
8. Search postings & \code{jobsearch/} & Postings, and companies for a spontaneous application, ranked by meaning. \\
9. Keep            & \code{store.py}, \code{companies.py} & Results saved at once; starred companies get a card on demand. \\
10. Apply          & \code{mail/}, \code{export/} & E-mail \emph{draft}; card or application as PDF, Excel, Word. \\
11. Follow         & \code{applications/} & What you sent, follow-ups, appointments. \\
12. Ask            & \code{chat/}, \code{mcp\_server.py} & Optional strip: your own Claude Code or Codex answers from your data. \\
\bottomrule
\end{tabularx}
}


\section{Design principles}
\begin{itemize}
  \item \textbf{Rules, not a chatbot.} Profile, scoring and affinity use rules, dictionaries and a
        small local embedding model; the program calls no large language model. Only \emph{your own}
        Claude Code or Codex is started, by the assistant strip (chapter~\ref{ch:mcp}) and, when you
        ask, to suggest companies (chapter~\ref{ch:interest}).
  \item \textbf{The CV stays on the PC.} It never goes to company websites; only a cleaned role title
        goes to the public ESCO API. The one exception is the \emph{anonymised summary} you see before
        the first suggestion request.
  \item \textbf{Nothing is sent for you.} Applications open as \emph{drafts}; you press Send (chapter~\ref{ch:apply}).
  \item \textbf{No long freezes, polite network.} Slow work runs in the background (section~\ref{sec:robust});
        pauses per host, robots.txt, caches; a site whose terms forbid automated reading is never contacted (chapter~\ref{ch:jobs}).
  \item \textbf{Nothing invented, estimates labelled.} Unknown values stay blank; simulators, the AI-writing
        figure, the affinity, a hidden employer and every suggestion say they are estimates or hypotheses.
\end{itemize}

\begin{infobox}
Scores are indicative: they help to decide where to look first, they are not a verdict.
\end{infobox}
